\chapter{Statistics from Repetition}
\label{chap:statistics-from-repetition}

Part III turns to heat, entropy, and matter --- the physics of the many. It opens not with a thermodynamic
law but with the statistics those laws rest on, because the instrument's count, gathered over repetition,
becomes a distribution before it becomes a temperature. The point worth holding onto is that the statistics
are \emph{earned}, not assumed: the instrument does not begin with a bell curve or a significance test and
apply it; it repeats a measurement, accumulates the counts, and reads the distribution off what accumulates.
Four readings carry the chapter --- the bell curve as the shape repetition makes, the two faces of the
disciplined small sample, and the posterior as the smallest correction a new fact can force --- and each is a
count gathered, not a law imposed.

\section{Gauss's First Effect: the bell curve as a stable shape}
\label{se:gauss-first}

Repeat a measurement many times and bin the outcomes, and the histogram settles into a shape. Not just any
shape: as the trials accumulate, the bins fill toward a single, stable, symmetric curve --- the bell ---
characterized completely by where it centers and how widely it spreads,
\begin{equation}
  p(x) = \frac{1}{\sqrt{2\pi\sigma^2}}\,\exp\!\left(-\frac{(x-\mu)^2}{2\sigma^2}\right).
\end{equation}
The instrument does not impose this curve; it watches it emerge. The bin totals add over concatenated trials
--- run the experiment twice as long and each bin's count roughly doubles, the shape unchanged --- and the
two parameters $\mu$ and $\sigma^2$ are read off the accumulated counts, the center and the width the data
themselves settle into. Gauss met the curve in the reduction of astronomical observations, where a star's
position, measured again and again, scattered about a central value in exactly this form, and he took it not
as a law of nature but as the stable shape of error.

Why this shape and no other? Because a single measurement's error is itself the sum of many small,
independent nudges --- a little turbulence in the air, a little tremor in the hand, a little flicker in the
apparatus --- and a sum of many small independent contributions, whatever each one's own distribution, piles
up into the bell. That is the reason the curve appears wherever a count is repeated: not because nature
prefers it, but because adding up independent small errors has only one stable shape to settle into.

That intuition is a theorem, the central limit theorem, and it is worth stating with its mechanism. Take $n$
independent contributions, each with a finite mean and variance but otherwise of any shape whatever, sum them,
and center and scale the sum by its own spread; as $n$ grows the distribution of that normalized sum converges
to the standard Gaussian, regardless of what the individual contributions looked like. The universality ---
for almost any underlying distribution --- is the whole point, and the reason it holds is visible in the
moment-generating function, the transform
\begin{equation}
  M_X(t) = \mathbb{E}\bigl[e^{tX}\bigr],
\end{equation}
which packages all of a distribution's moments into one function. Its defining virtue is that the
moment-generating function of a sum is the product of the moment-generating functions: independent
contributions multiply their transforms, so their logarithms add. Take logarithms, then, and normalize, and
every term in the log-transform beyond the second is suppressed by a growing power of $n$ and dies, leaving
only a quadratic. A quadratic log-transform is exactly the signature of the Gaussian, and nothing else. So the
bell is not one stable shape among several; it is the unique fixed point of adding-and-normalizing, the only
distribution whose log-transform a sum cannot push away from. The instrument does not impose the bell; it is
the only shape a sum of many small independent counts can settle into.

The
honest floor is a finite count obligation: the normal curve is the stable shape of the bin totals, an
emergent count and not a posited law. The reading is that the bell curve is not a fact about nature but the
shape that repetition makes --- what a sum of many independent finite errors converges to, characterized by
its center and its spread and needing nothing more.

\section{The Gosset Effect, I: the disciplined small sample}
\label{se:gosset-pooling}

The bell curve is the shape of \emph{many} repetitions. But the instrument is often not given many --- it has
a handful of measurements and must say, honestly, how much it knows from them. This is the first of Gosset's
two problems, the problem of pooling. With only a few samples the true spread of the underlying distribution
is unknown, so the instrument cannot reach for the population's variance; it must estimate the spread from the
same small sample it is trying to measure. The mean it forms is the signal, and the signal sharpens as the
sample grows: the random errors in the individual measurements partly cancel when averaged, so the noise in
the mean pools down as $\sqrt{n}$, and the honest interval around the mean is
\begin{equation}
  \bar{x} \pm t\,\frac{s}{\sqrt{n}},
\end{equation}
the sample mean give or take a few standard errors $s/\sqrt{n}$. The multiplier $t$ is the crucial honesty.
It is drawn not from the normal curve but from Student's $t$-distribution, which has heavier tails --- and the
heavier tails are exactly the price of estimating the spread from the small sample rather than knowing it.
With few measurements the instrument is unsure not only of the mean but of its own estimate of the scatter,
and the fatter tails of the $t$ widen the interval to pay for that second uncertainty rather than pretend it
away.

The shape of those tails has a precise form. The $t$-density, with $\nu = n - 1$ degrees of freedom, falls off
as a power of the deviation,
\begin{equation}
  p(t) \propto \left(1 + \frac{t^2}{\nu}\right)^{-(\nu+1)/2},
\end{equation}
rather than the Gaussian's exponential $e^{-t^2/2}$: a polynomial tail, not an exponential one, so extreme
deviations are far likelier under the $t$ than under the bell. With few degrees of freedom the tails are
heavy and the interval wide; as $\nu$ grows --- as the sample fills out and the estimated spread approaches the
true one --- the power-law tail steepens toward the Gaussian's, and in the limit $\nu \to \infty$ the $t$
becomes the bell exactly. The $t$-distribution is the bell with its own uncertainty about its width built
honestly into the tails, and it relaxes to the bell precisely as that uncertainty is resolved.

The effect is named for ``Student,'' and the pseudonym is part of the lesson. William Gosset discovered the
distribution while working as a brewer at Guinness, where the samples were necessarily small --- a few batches
of barley, a handful of yeast cultures --- and the brewery forbade its staff from publishing under their own
names. The $t$-distribution is, at bottom, the honest statistics of the under-sampled: what an instrument may
claim when it has only a little data and refuses to borrow confidence it has not earned. The honest floor is a
finite count obligation: the pooled estimate is a count of how fast the noise in the mean shrinks with $n$,
and the $t$ multiplier is the honest widening that a small count demands.

\section{The Gosset Effect, II: hypothesis as projection}
\label{se:gosset-projection}

Gosset's second problem is different in kind, and it deserves its own room: not how much the sample says, but
whether an apparent signal in it is real. A sample mean by itself means nothing until it is set against the
variability that could have produced it by chance, and the instrument settles the question by a projection.
Take the data as a vector $x$, and let a hypothesis name an axis $u$ --- the direction the data would point if
the hypothesis held. Project $x$ onto $u$. The component along the axis, $\langle x, u\rangle$, is the
structure the hypothesis predicts; the part left over, $x - \langle x, u\rangle\,u$, lies orthogonal to $u$
and is everything the hypothesis does not explain. The test statistic is the ratio of the two,
\begin{equation}
  t = \frac{\langle x, u\rangle}{\lVert x - \langle x, u\rangle\, u\rVert},
\end{equation}
the structure along the axis measured against the residual that scatters away from it: large when the data
line up with the hypothesis, small when they spill into the orthogonal directions the hypothesis has nothing
to say about.

This is the same geometry the kernel chapter ran on a wave. There the instrument projected a record onto a
subspace and set aside the orthogonal residual it could not reach; here it projects the data onto the
hypothesis axis and reads the structure against the orthogonal residual that could be noise. A significance
test is a Hilbertian projection --- a structured comparison of an observation against a hypothesis, decided by
how much of the observation lies along the hypothesis and how much falls away from it. The honest floor is a
smooth-shadow analogy: the discrete projection-and-residual is the shadow of the continuum $t$-distribution,
and the bridge between the instrument's finite decomposition and the smooth curve of critical values is an
analogy, named as one.

The two readings together are the discipline Gosset brought to a small sample. The first asks how much the
sample knows and widens the answer honestly for the count it has; the second asks whether a pattern in the
sample is real and settles it by projecting onto the hypothesis. Neither helps itself to the distribution it
is trying to test; both are bookkeeping the instrument can do on a handful of counts, refusing to borrow what
it has not measured.

\section{The Bayes Effect: the minimal coherence-restoring correction}
\label{se:bayes}

The last reading is about how the instrument changes its mind. A new fact does not overwrite what was known;
it corrects it by the least amount that restores coherence. The instrument carries a prior --- its accumulated
state of belief --- and a new event brings a likelihood, the constraint that event imposes; the posterior is
the unique reconciliation that minimizes the strain between the two, the smallest revision that makes the old
belief and the new fact consistent. In logarithmic form the update is simply additive,
\begin{equation}
  \log P(\theta \mid D) = \log P(\theta) + \log P(D \mid \theta) + c,
\end{equation}
the log-posterior being the log-prior plus the log-likelihood, up to a normalizing constant $c$. Read
descriptively, this says that evidence accumulates by addition: the instrument does not throw away the prior
and start over; it adds the new log-evidence to what it already held and renormalizes, and the result is the
least correction that restores coherence.

The honest floor is a finite count obligation --- the update is an additive count of log-evidence, the
minimal correction that restores coherence. The reading is that the posterior is not a fresh belief replacing
an old one but the smallest correction to the old that the new fact forces. What was known is carried forward,
corrected rather than erased, and the cost of the correction is exactly the surprise the new fact carried ---
the log-evidence it added. The instrument changes its mind the way it does everything else: by the least
bookkeeping that keeps the account consistent.

\bigskip

With these readings the instrument's statistics are complete, and they were earned, not assumed. A count
repeated settles into the bell curve, the shape any sum of small independent errors must take. A small count
is read with discipline --- pooled honestly for what it knows, projected to test what it claims. And a new
fact revises the old by the smallest correction that restores coherence. Nowhere did the instrument begin with
a law of statistics and apply it; everywhere it gathered counts and read the distribution off them. Part III
turns next from the distribution to what a distribution costs --- to entropy, the second law, and the price
the instrument pays to know.
